\section{Заключение}
\label{sec:Chapter5} \index{Chapter5}

В рамках выпускной квалификационной работы исследована и реализована
нейросетевая система восстановления и очистки речевых сигналов, основанная на
архитектуре Miipher-2 и самоконтролируемых акустических представлениях. Главной
решённой проблемой стала недоступность ключевого компонента оригинальной модели
--- закрытого энкодера USM --- и его замена открытыми аналогами.

В ходе работы получены следующие основные результаты.
\begin{enumerate}
    \item Выполнен аналитический обзор методов улучшения и восстановления речи
        --- от классической спектральной обработки до генеративных и
        SSL-подходов, --- а также SSL-энкодеров и нейросетевых вокодеров; на его
        основе обоснован выбор схемы Miipher-2 как наиболее соответствующей
        требованиям.
    \item Формализована задача восстановления речи в пространстве признаков
        SSL-энкодера и сформулированы функциональные, качественные и
        вычислительные требования к решению.
    \item Воспроизведена архитектура Miipher-2 на полностью открытых
        компонентах: энкодерах WavLM Base+ и HuBERT Base, очистителе на основе
        параллельных адаптеров и вокодерах HiFi-GAN и WaveFit.
    \item Экспериментально установлено, что качество восстановления немонотонно
        зависит от номера слоя извлечения признаков и достигает максимума на
        средних слоях ($\ell=7$ для WavLM, $\ell=6$ для HuBERT), что согласуется
        с известными данными о структуре SSL-представлений.
    \item Проведено сравнение четырёх комбинаций <<энкодер — вокодер>> по
        метрикам PESQ, STOI, SI-SDR и DNSMOS, а также по вычислительной
        эффективности. Все конфигурации превзошли целевые ориентиры по приросту
        качества.
    \item Определена рекомендуемая конфигурация \textbf{WavLM Base+ + HiFi-GAN}
        (слой~7), обеспечивающая наилучший баланс качества и быстродействия
        (PESQ~$2{,}94$, STOI~$0{,}927$, $\mathrm{RTF}=0{,}034$).
\end{enumerate}

Таким образом, поставленная цель достигнута: показано, что закрытый энкодер USM
может быть заменён открытыми SSL-моделями без существенной потери качества, а
выбор энкодера и слоя извлечения признаков влияет на результат сильнее, чем тип
вокодера. Практическая значимость работы состоит в возможности построения
полностью открытой системы очистки речи для автоматического формирования
обучающих корпусов студийного качества.

В качестве направлений дальнейшего развития можно выделить: применение более
крупного энкодера WavLM Large и анализ его слоёв; совместное (end-to-end)
дообучение очистителя и вокодера; расширение на многоязычные данные; а также
субъективную оценку качества методом MOS с привлечением аудиторов.
